% !TeX root = ../main.tex
\chapter{Bozze e Argomenti in Sospeso}

\section{Argomenti in sospeso da riprendere}

\textbf{Cap.4}:


\begin{itemize}
    \item \textbf{Factory Methods}: par. 4.4.4.
    \item \textbf{Distruzione di oggetti}: par. 4.6.8.
    \item \textbf{Costruttori canonici nei records}: par. 4.7.2.
    \item \textbf{seconda parte delle lambdas da fare/rivedere}: a partire da par. 6.2.5.
\end{itemize}



\section{Il Class Path}

Il \textbf{Class Path} (percorso delle classi) è la collezione di tutte le locazioni (cartelle o archivi) in cui la JVM e il compilatore cercano i file \texttt{.class}. È, di fatto, l'elenco dei punti di partenza per la risoluzione dei package.

Il Class Path è una stringa composta da vari percorsi. Il separatore dipende dal sistema operativo:

\begin{table}[h]
    \centering
    \small % Riduce leggermente la dimensione del testo per far stare tutto comodamente
    \begin{tabularx}{\textwidth}{l X l} % Rimosse le linee verticali per stile professionale
        \toprule
        \textbf{Sistema} & \textbf{Esempio di Class Path} & \textbf{Separatore} \\
        \midrule
        \textbf{Fedora / UNIX} & \path{/home/user/classdir:.:/libs/archive.jar} & Due punti (\texttt{:}) \\
        \addlinespace % Aggiunge un piccolo spazio tra le righe
        \textbf{Windows} & \path{C:\classdir;.;C:\libs\archive.jar} & Punto e virgola (\texttt{;}) \\
        \bottomrule
    \end{tabularx}
\end{table}

\begin{observation}[Il simbolo del punto]
    In entrambi i sistemi, il punto (\texttt{.}) rappresenta la \textbf{directory corrente}.
\end{observation}

È comune confondere il Classpath con la directory di output. Ecco la distinzione corretta:

\begin{itemize}
    \item \textbf{Per il Compilatore (\texttt{javac})}: Il Classpath serve a \textbf{risolvere i riferimenti esterni}. Se la tua classe \texttt{A} usa la classe \texttt{B} (che si trova in un JAR o in un'altra cartella), il compilatore deve "leggere" \texttt{B} per verificare che tu stia usando i suoi metodi correttamente.
    \begin{itemize}
        \item \textit{Nota}: Per decidere \textbf{dove mettere} il file compilato, si usa sempre il flag \texttt{-d}, non il Classpath.
    \end{itemize}

    \item \textbf{Per la JVM (\texttt{java})}: Il Classpath serve a \textbf{caricare il codice in memoria}. La JVM non sa nulla dei file \texttt{.java}; lei cerca solo il bytecode (\texttt{.class}) seguendo i percorsi indicati nel Classpath per poter avviare l'esecuzione.
\end{itemize}

\subsection{La "Ricerca Intelligente" del Compilatore}
Horstmann sottolinea che il compilatore è più sofisticato della JVM nella ricerca. Quando incontra una classe, ad esempio \texttt{Employee}:
\begin{enumerate}
    \item Cerca il file \texttt{Employee.class} nel Classpath.
    \item Se trova anche il file \texttt{Employee.java}, controlla le date di modifica.
    \item \textbf{Auto-recompile}: Se il sorgente (\texttt{.java}) è più recente del compilato (\texttt{.class}), il compilatore ricompila automaticamente il sorgente per garantire che tu stia usando la versione più aggiornata.
\end{enumerate}

\subsection{I file JAR (Java Archive)}
Oltre alle normali directory, le classi possono essere raggruppate in file \textbf{JAR}.
\begin{itemize}
    \item Un file JAR è essenzialmente un archivio in formato \textbf{ZIP} che contiene classi e strutture di sottocartelle compresse.
    \item È lo standard per la distribuzione di librerie di terze parti.
    \item Permette di migliorare le prestazioni (un solo file da aprire invece di centinaia) e risparmiare spazio.
\end{itemize}

\subsection{L'Algoritmo di Ricerca}
Se la JVM deve caricare la classe \texttt{com.horstmann.corejava.Employee}, cercherà in ordine:
\begin{enumerate}
    \item Nelle classi delle \textbf{API di Java} (sempre cercate per prime, non serve includerle nel Class Path).
    \item In ogni voce del Class Path, aggiungendo il percorso del package alla base. Ad esempio, se il path contiene \texttt{/home/user/classdir}, cercherà: \\
    \texttt{/home/user/classdir/com/horstmann/corejava/Employee.class}
\end{enumerate}

\begin{caution}[La trappola del punto]
    Per impostazione predefinita, il Class Path è costituito solo dal punto (\texttt{.}). Tuttavia, se specifichi un Class Path personalizzato (tramite la variabile d'ambiente \texttt{CLASSPATH} o il flag \texttt{-cp}) e \textbf{dimentichi} di includere il punto (\texttt{.}), la JVM non guarderà più nella cartella corrente.
    \begin{itemize}
        \item \textbf{Risultato}: Il programma compila (perché \texttt{javac} guarda sempre la cartella corrente), ma non parte (perché \texttt{java} non la trova).
    \end{itemize}
\end{caution}

\subsection{Uso delle Wildcard}
È possibile includere tutti i file JAR di una cartella usando l'asterisco:
\begin{lstlisting}[language=bash]
# Su Fedora (bisogna proteggere l'asterisco dalle espansioni della shell)
java -cp "/home/user/libs/*:." MiaClasse
\end{lstlisting}
\textbf{Nota}: L'asterisco include tutti i file \texttt{.jar}, ma non i singoli file \texttt{.class} sparsi nella cartella.

\section{bozze eccezioni}

\begin{tcolorbox}[colback=blue!5!white,colframe=blue!75!black,title=Best Practice: Eccezioni vs Performance]
\textbf{1. Usa le eccezioni solo dove necessario}: Non utilizzare mai le eccezioni per gestire il normale flusso di controllo del programma (es. per uscire da un loop). Creare un'eccezione è "costoso" per la JVM perché deve catturare l'intero \textit{stack trace} (lo stato di tutte le chiamate ai metodi in quel momento).

\textbf{2. Preferisci metodi che lanciano eccezioni}: Quando esplori le API di Java, preferisci sempre i metodi che segnalano errori tramite eccezioni rispetto a quelli che restituiscono codici di errore (come \texttt{null} o \texttt{-1}).
\begin{itemize}
\item Un'eccezione \textbf{non può essere ignorata}: se non la gestisci, il programma si ferma, impedendo al bug di propagarsi silenziosamente.
\item Un codice di errore richiede un \texttt{if} manuale: se il programmatore si dimentica di scriverlo, il programma continuerà a girare in uno stato instabile.
\end{itemize}
\end{tcolorbox}

\begin{tcolorbox}[colback=blue!5!white,colframe=blue!75!black,title=Best Practice]
Affidarsi esclusivamente al \textit{Default Exception Handler} della JVM come strategia ordinaria di gestione degli errori è generalmente una cattiva pratica. Una progettazione solida deve mirare alla robustezza e alla resilienza del software.

Le eccezioni dovrebbero essere intercettate nel livello più basso che sia effettivamente in grado di gestirle in modo corretto; in caso contrario, è preferibile propagarle verso livelli superiori che dispongano del contesto necessario per decidere come reagire.
\end{tcolorbox}

\subsection{Checked Exceptions}
Da quanto detto finora, sono \textit{Checked Exceptions} tutte le \texttt{Exception} che non siano figlie di \texttt{RunTimeException}; abbiamo anche detto che il compilatore ci obbliga a \textit{gestire} questo tipo di eccezioni: vediamo cosa vuol dire.  

Il compilatore Java agisce come un guardiano rigoroso per le \textit{Checked Exceptions}, applicando la regola del \textbf{Catch or Declare}:

\begin{tcolorbox}[colback=blue!5!white,colframe=blue!75!black,title=Regola del Catch-or-Declare]
Se un metodo X ne invoca un altro che lancia una checked exception, il codice \textbf{non compila} finché non si soddisfa una di queste due condizioni:
\begin{enumerate}
    \item \textbf{Catch}: Si circonda l'istruzione con un blocco \texttt{try-catch} per gestire l'errore.
    \item \textbf{Declare}: Si aggiunge la clausola \texttt{throws} alla firma del metodo X per delegare la responsabilità a chiunque lo invochi.
\end{enumerate}
\end{tcolorbox}

Le \textbf{Unchecked Exceptions} ignorano questa regola: il compilatore assume che tu scriva codice corretto per evitarle (ad esempio controllando che un oggetto non sia \texttt{null} prima di usarlo), quindi non ti obbliga a usare blocchi \texttt{try-catch}.

\subsection{Dichiarare e Lanciare Eccezioni}

Un metodo Java può terminare in due modi: restituendo un valore o "esplodendo" con un'eccezione. Se un metodo non è in grado di gestire una situazione anomala, deve dichiarare questa possibilità al mondo esterno. La sintassi per dichiarare e \textit{lanciare} un'eccezione è illustrata nel seguente esempio elementare: 

\begin{lstlisting}[caption={Esempio di metodo che lancia un'eccezione}]
public void setAge(int age) throws IllegalArgumentException {
    if (age < 0) {
        // Creiamo l'oggetto eccezione e lo lanciamo
        throw new IllegalArgumentException("age cannot be negative: " + age);
    }
    this.age = age;
}
\end{lstlisting}

È fondamentale non confondere queste due parole chiave, simili ma con scopi diversi:
\begin{itemize}
    \item \textbf{\texttt{throws}} (al plurale): Si usa nella \textbf{firma del metodo} per avvertire il compilatore e il chiamante che quel metodo può lanciare determinate eccezioni.
    \item \textbf{\texttt{throw}} (al singolare): Si usa nel \textbf{corpo del metodo} per lanciare effettivamente l'oggetto eccezione.
\end{itemize}


\begin{observation}
Cosa succede quando viene eseguita l'istruzione \texttt{throw}?
\begin{enumerate}
    \item L'esecuzione del metodo si interrompe \textbf{immediatamente}.
    \item Le istruzioni successive (come l'assegnamento \texttt{this.eta = eta}) non vengono eseguite.
    \item Il controllo passa al chiamante, cercando il primo blocco \texttt{try-catch} disponibile che sappia gestire \texttt{IllegalArgumentException}.
\end{enumerate}
\end{observation}

\begin{tip}
In questo esempio abbiamo usato \texttt{IllegalArgumentException}, che è una \textbf{Unchecked Exception}. Per questo motivo, la clausola \texttt{throws} nella firma del metodo sarebbe tecnicamente opzionale per il compilatore, ma è considerata un'ottima pratica di documentazione per avvisare chi userà il tuo metodo.
\end{tip}

\subsection{Quando usare la clausola \texttt{throws}}

Non è necessario dichiarare ogni possibile oggetto \texttt{Throwable} che il tuo metodo potrebbe lanciare. Per capire cosa va inserito nella clausola \texttt{throws}, dobbiamo distinguere quattro situazioni in cui può verificarsi un'eccezione:

\begin{enumerate}
    \item \textbf{Chiamata a un metodo che lancia una checked exception}: ad esempio, l'apertura di un file tramite il costruttore di \texttt{FileInputStream}.
    \item \textbf{Rilevamento di un errore esplicito}: quando decidi di lanciare una checked exception usando l'istruzione \texttt{throw}.
    \item \textbf{Errore di programmazione (Bug)}: ad esempio un accesso a un indice errato \texttt{a[-1] = 0} (genera una \texttt{ArrayIndexOutOfBoundsException} - \textit{unchecked}).
    \item \textbf{Errore interno della JVM}: ad esempio l'esaurimento della memoria (\texttt{Error} - \textit{unchecked}).
\end{enumerate}

\begin{important}
Devi dichiarare nella firma del metodo (tramite \texttt{throws}) solo le eccezioni appartenenti ai primi due scenari (\textbf{Checked Exceptions}).
\end{important}



\subsection{Metodi come "Trappole Mortali"}

Horstmann definisce i metodi che lanciano eccezioni non dichiarate come potenziali "trappole mortali". Se un'eccezione viene lanciata e nessun \textit{handler} la cattura, il thread di esecuzione corrente termina immediatamente.

\begin{lstlisting}[caption={Esempio di propagazione (Death Trap)}]
// Poiche' FileInputStream lancia una FileNotFoundException (checked),
// siamo OBBLIGATI a dichiararla con 'throws'
public void leggiFile(String nome) throws FileNotFoundException {
    FileInputStream f = new FileInputStream(nome);
    // ... logica di lettura ...
}
\end{lstlisting}

\begin{warning}
Se ometti la clausola \texttt{throws} in presenza di una checked exception, il compilatore Java genererà un errore. Al contrario, non dovresti mai dichiarare le \texttt{RuntimeException} (bug tuoi) o gli \texttt{Error} (guasti JVM), poiché sono fuori dal controllo della normale logica di gestione degli errori.
\end{warning}

